\noindent\textbf{Background:} Two-dimensional maps of trial-averaged responses, made with t-SNE or
multidimensional scaling, are read neighbor by neighbor, and analysts check them against the
averaged data. Because the map is fitted to the same noisy averages, such checks share its errors.

\noindent\textbf{New method:} We define the false comfort of a check as the share of the
neighborhoods it passes that are not faithful to the noise-free responses. We estimate it from
held-out repetitions, with a simulation--extrapolation correction for the reference's own noise and
subject-bootstrap bounds, and validate the estimator in plasmode simulations.

\noindent\textbf{Results:} Pre-specified analyses used human EEG (THINGS-EEG2; 10 subjects,
200 images) and two-photon imaging of mouse V1 (Allen Brain Observatory; 10 mice, 118 images). At 8 trials, \EegFCeight\% of the neighborhoods that passed the plug-in check were not
faithful in EEG (conservative bound \EegLeight\%), and essentially all were in mouse V1 (bound
\MouseLeight\%). Maps of cross-validated distances passed fewer neighborhoods, most of them
unfaithful. The result held in \EegCells{} and \MouseCells{} of 12 pre-specified variations.

\noindent\textbf{Comparison with existing methods:} Of six common checks computed from the
analyst's own trials, per-point trustworthiness, scDEED, cross-validated distances and recurrence
across t-SNE seeds left most false comfort in both datasets. Bootstrap and split-half recurrence
passed almost no neighborhoods at 8 trials. None removed most false comfort.

\noindent\textbf{Conclusions:} Checks computed from the same trials cannot reveal where trial
averages are wrong. Local neighborhoods of such maps should be read against held-out repetitions,
which the estimator makes practical.
